% =====================================================================
%  COVER LETTER -- submitted as a separate file.
%  Every figure quoted here is taken from the manuscript and must stay
%  identical to it.
% =====================================================================
\documentclass[11pt]{article}
\usepackage[a4paper,hmargin=2.2cm,vmargin=1.6cm]{geometry}
\usepackage{times}
\pagestyle{empty}
\setlength{\parindent}{0pt}
\setlength{\parskip}{0.45em}

% British date format, to match the manuscript's spelling.
\renewcommand{\today}{\number\day\space\ifcase\month\or January\or February\or
  March\or April\or May\or June\or July\or August\or September\or October\or
  November\or December\fi\space\number\year}

\begin{document}

\begin{flushright}
Indrajeet Mondal\\
Department of Computer Application, Narula Institute of Technology\\
81 Nilgunj Road, Agarpara, Kolkata 700109, India\\
indrajeetmondal120703@gmail.com\\[0.4em]
\today
\end{flushright}

The Editor-in-Chief\\
\emph{Future Generation Computer Systems}

{\raggedright\textbf{Submission of a Full Length Article: ``The escalation tax:
measuring what manual model and reasoning-mode selection costs, and whether a
one-shot router recovers it''}\par}

Dear Editor,

We submit the above manuscript for consideration as a Full Length Article in
\emph{Future Generation Computer Systems}.

Language-model services are now sold in tiers, and a growing number expose a
second control beside model size: whether the model should reason before it
answers. Users allocate by trial and error: some start small and move up after
every wrong answer, paying for each discarded attempt, while others leave a
large reasoning configuration selected and pay its premium on prompts that
never needed it. None of the routing work we located measures what this manual
behaviour costs, or charges a router for its own energy.

We define the \emph{escalation tax}, the cost a policy pays beyond the cheapest
configuration that would have answered, and split it into escalation and
stickiness. We measure it on eight Qwen3 configurations (0.6B to 8B parameters,
reasoning on and off), running each of 1{,}200 gradable prompts through every
configuration on a laptop GPU with metered board power, and on two further
pools: thirteen public models priced by their providers, and six hosted
configurations with paid reasoning effort. Classical, encoder and
language-model routers are compared with manual policies and a self-verifying
cascade, net of each router's own tokens, latency and energy. The main findings
are:
\begin{itemize}\setlength{\itemsep}{0.15em}
  \item Leaving the largest reasoning configuration selected pays 12.96 times
        the minimum sufficient energy, while a frugal climb that moves up in
        size before turning reasoning on pays 0.89 times. The same ordering
        holds on the public pool (7.96 against 2.69) and the hosted pool (4.35
        against 0.22).
  \item Below 8B, a larger model buys accuracy for 3.8 to 5.5 times less energy
        per solved prompt than turning reasoning on. Energy per token rises
        6.6-fold from 0.6B to 8B, which token counts cannot see.
  \item Whether a one-shot router beats the frugal climb depends on how a failed
        route is recovered. Resent to the largest configuration, the best router
        spends 1.30 times the climb's energy per solved prompt, though 64\% less
        than the sticky default; recovered by climbing, every router comes
        within 6\% of the climb and four are cheaper. Where users miss wrong
        answers, the router answers 7.1 points more.
\end{itemize}

The work suits \emph{Future Generation Computer Systems} because serving
language models across configurations of differing size, reasoning effort and
price is a resource-allocation problem for computing infrastructure, within the
journal's scope in distributed dynamic resource management and scheduling,
internet computing, and scaling and performance. It contributes a cost model over heterogeneous
configurations, measured rather than estimated energy, and an evaluation
spanning local hardware, public models and a hosted service, with concrete
guidance for operators: open on a small configuration with reasoning off, and
recover a failed route by climbing rather than jumping to the largest one.

The manuscript is original, has not been published, and is not under
consideration elsewhere. Both authors have approved it and agree to its
submission. We have no competing interests, the work received no specific
funding, and the use of generative AI is declared in the manuscript. It keeps
to the 18-page limit for Full Length Articles; supplementary material
accompanies it. Thank you for considering our submission.

Yours sincerely,

Indrajeet Mondal\\
Corresponding author, on behalf of both authors\\
ORCID: 0009-0009-7966-2683

\end{document}
